\documentclass[11pt]{article}
\usepackage{amsmath,amssymb,amsthm,hyperref,geometry,longtable}\geometry{margin=1in}
\newtheorem{theorem}{Theorem}\newtheorem{lemma}{Lemma}
\title{A display-exact native interaction criterion}
\date{15 November 2026}\begin{document}\maketitle
\paragraph{Status and order of construction.} General categorical theorem proposed for independent review. The context and content inputs below are specified independently; the crossing and its comparisons are derived, not supplied. Semantic realizations are not prerequisites to this proof. The criterion is complementary to the existing rig-based O1 theorem: it permits an arbitrary pullback base and Cartesian occurrence monad, but restricts contexts to actual display maps. It is not a strict generalization of every nonfaithful rig-coefficient context, and calling a display ``joint'' alone does not certify an arbitrary physical interference model.

\section{Independent data}
Let $\mathcal C$ be a locally small category with chosen pullbacks in suitable nested universes. Let $\mathcal D$ be a class of arrows containing identities and isomorphisms and closed under pullback and composition. The context fibre $K_A$ is the full subcategory of $\mathcal C/A$ on arrows $\pi:E\to A$ in $\mathcal D$. Its extension is $E$, its pure reindexing is pullback, its unit is $\mathrm{id}_A$, and its dependent accumulation is composition of display maps. These operations give the display-map contextad; its native arrows are all commuting slice arrows, not only isomorphisms.

Independently choose a Cartesian monad $(S,\eta^S,\mu^S)$ on $\mathcal C$, with $S(\mathcal D)\subseteq\mathcal D$. Here Cartesian means that $S$ preserves pullbacks and its unit and multiplication are Cartesian natural transformations. No polynomial representation, finite positions or tensor of independent effects is required for this theorem.

Choose an independently valid contentad with projection $p:L\to\mathcal C$, extension $d:L\to\mathcal C$, unit $U$, typed accumulation $a(Q,R)=Q\diamond R$ for $pR=dQ$, and actual extension bodies $\eta^L_A:A\to dU_A$, $\mu^L_{Q,R}:dR\to d(Q\diamond R)$. In particular its one-sided relative pseudopromonad and global opcartesian transport already satisfy their own laws. Do not infer a contentad merely from an opfibration or a slice factorization.

The additional assumptions concern this content input and the fixed display class:
\begin{enumerate}
\item $p$ has genuine chosen Cartesian lifts $c_{f,Q}:f^*Q\to Q$ for $f\in\mathcal D$, besides its global opcartesian lifts.
\item A natural map $\nu:d\Rightarrow Sp$ satisfies
\[
 \nu_{U_A}\eta^L_A=\eta^S_A,\qquad
 \nu_{Q\diamond R}\mu^L_{Q,R}=\mu^S_A\,S\nu_Q\,\nu_R.
\tag{V}
\]
\item For every $f:X\to A$ in $\mathcal D$, the natural square
\[
 \begin{array}{ccc}d(f^*Q)&\xrightarrow{dc_{f,Q}}&dQ\\
 \nu_{f^*Q}\downarrow&&\downarrow\nu_Q\\
 SX&\xrightarrow{Sf}&SA\end{array}
\tag{D}
\]
is a pullback.
\item The unit $U$ preserves these Cartesian arrows. Typed accumulation preserves the Cartesian pair $(c_{f,Q},c_{j_f,R})$, where $j_f=dc_{f,Q}$. This is well-typed because (D) makes $j_f$ a pullback of $Sf$, hence $j_f\in\mathcal D$.
\item For every pullback square with vertical arrows in $\mathcal D$, the canonical mate between Cartesian pullback and opcartesian transport is invertible. For
$\pi_u:E_u\to A$, $\pi':E'\to A'$ and horizontal maps $v:E_u\to E'$, $u:A\to A'$, it is
\[
 v_!\pi_u^*Q\ \cong\ \pi'^*u_!Q.
\tag{BC}
\]
This is the canonical factorization mate, not a freely supplied mixed coherence map.
\end{enumerate}
These are abstract counterparts of O1's content-only display, Cartesian accumulation and Beck--Chevalley hypotheses. Their validity is a substantive input. No context/content crossing occurs in their statement.

\section{Constructed crossing and all native arrows}
For $\pi:E\to A$ and $Q\in L_A$, put
\[
 H_\pi(Q)=dQ\times_{SA}SE,\qquad
 \widehat\pi:H_\pi(Q)\to dQ,\qquad \widehat Q=\pi^*Q.
\tag{X}
\]
Since $S\pi\in\mathcal D$, $\widehat\pi$ is an admitted context. Display (D) gives a canonical isomorphism
$\chi_{\pi,Q}:H_\pi(Q)\cong d(\pi^*Q)$ preserving both projections, explicitly
\[d(c_{\pi,Q})\chi_{\pi,Q}=\operatorname{pr}_1,\qquad \nu_{\pi^*Q}\chi_{\pi,Q}=\operatorname{pr}_2.\]
For a general represented content-then-context operation, with $f:B\to dQ$ and $g:E\to C$, pull back $\widehat\pi$ along $f$, push $\widehat Q$ along $g$, and take the body
\[
 \operatorname{dom}(f^*\widehat\pi)\to H_\pi(Q)
 \xrightarrow{\chi_{\pi,Q}}d(\pi^*Q)
 \xrightarrow{d\kappa_g}d(g_!\pi^*Q).
\tag{G}
\]

For simultaneous pure/native maps $u:A\to A'$, $h:E\to E'$ with $\pi'h=u\pi$, and $b:Q\to Q'$ over $u$, the context map in (X) is the unique pullback map with coordinates $db$ and $Sh$. The content map $\pi^*Q\to\pi'^*Q'$ over $h$ is the Cartesian factorization of $b c_{\pi,Q}$ through $c_{\pi',Q'}$. Identities and compositions of grade maps follow from Cartesian uniqueness. The body square follows from the two display projections. Thus all native and pure squares are retained. Equality of content grade maps is never inferred merely from equality after $d$.

For pure-middle balance along $u$, preservation of the context-reindexing pullback by $S$ gives the required context carrier isomorphism. The content component is (BC). Its mate equation and the target display pullback prove the body square. Naturality and identity/composition of balance follow from the defining universal factorizations and mate pasting. Hence (G) descends through the product arrows and pure-middle balance generators of the relative pseudocoend.

\section{The four structural comparisons}
\paragraph{Context unit.} For $\pi=\mathrm{id}_A$, (X) is the identity context on $dQ$, while $\pi^*Q\cong Q$. These are the canonical pullback/lift comparisons.

\paragraph{Content unit.} Restrict (X) for $U_A$ along the actual body $\eta^L_A$. By (V), its carrier is $A\times_{SA}SE$. Cartesianness of $\eta^S$ identifies it with $E$. The grade comparison $\pi^*U_A\cong U_E$ follows from preservation of Cartesian arrows by $U$. The body equation follows from these Cartesian factorizations and (V); no invertibility of $\eta^L$ is assumed.

\paragraph{Context multiplication.} For composable displays $F\xrightarrow\rho E\xrightarrow\pi A$, successive crossing has carrier
\[
 (dQ\times_{SA}SE)\times_{SE}SF\ \cong\ dQ\times_{SA}SF.
\]
This is the carrier of the single crossing for $\pi\rho$. The grade comparison is $\rho^*\pi^*Q\cong(\pi\rho)^*Q$. Pullback pasting and Cartesian lift composition give the native isomorphism and body equation.

\paragraph{Content multiplication.} Let $Q\in L_A$, $R\in L_{dQ}$, and $H=dQ\times_{SA}SE$. The successive-crossing carrier, before the final comparison, is
\[
 dR\times_{S(dQ)}SH
 \ \cong\ dR\times_{S^2A}S^2E.
\tag{M1}
\]
Here the map from $dR$ to $S^2A$ is $S\nu_Q\nu_R$, and the first isomorphism uses preservation of pullbacks by $S$. Cartesianness of $\mu^S$ then identifies (M1) with
\[
 dR\times_{SA}SE,
\tag{M2}
\]
where the map to $SA$ is $\mu^S_A S\nu_Q\nu_R$. By (V), (M2) is the direct crossing for $Q\diamond R$ restricted along its actual body $\mu^L_{Q,R}$. The output grade isomorphism is separately supplied by Cartesian preservation of typed accumulation:
\[
 \pi^*Q\diamond j_\pi^*R\ \cong\ \pi^*(Q\diamond R),
\]
using the display identification $H\cong d(\pi^*Q)$. Its defining Cartesian arrow and (V) give the comma body equation: the first display projection agrees by naturality of the actual body $\mu^L$ along the Cartesian pair, and the second agrees by (V) and naturality of $\mu^S$. No accumulated content body is identified with an entire extension by an assumed inverse.

\section{Full coherence, rather than an object formula}
\begin{theorem}[Display-exact interaction]
The independent data above determine a full native pseudodistributive law
$\tau:N\odot M\Rightarrow M\odot N$ for
\[
 M(A,B)=(e_A\downarrow B),\qquad N(A,B)=(A\downarrow d_B).
\]
It induces the relative pseudopromonad on the literal comma family $(e_A\downarrow d_B)$, retaining native noninvertible arrows and arbitrary pure endpoint squares.
\end{theorem}
\begin{proof}
Generator-level descent and the four invertible comparisons were constructed above. The remaining equations can be checked without element arguments: context objects are actual display maps, so their arrows are determined by the relevant pullback projections. Content comparison arrows are determined separately by composing with their fixed target Cartesian arrow. Under these tests the following list gives the common maps on both routes, using the original O1/Gambino--Lobbia C1--C10 convention.
\begin{description}
\item[C1, content associativity.] Expand to the pullback over $S^3A$ and $S^3E$. Both routes use the same composite $S^3E\to SE$ by associativity of $\mu^S$. The grade routes compose with the target Cartesian arrow to the same one-sided content associator.
\item[C2, content left unit.] On the representative restricted along the unit body, the two maps agree by $\mu^S(\eta^SS)=\mathrm{id}$ and the left unit of the contentad.
\item[C3, content multiplication against context unit.] Set $\pi=\mathrm{id}$ in (M1)--(M2). All display projections are identities; both routes give the same content multiplication body and its unit context.
\item[C4, both units.] Both maps are the same identity-display comparison over the pure unit body, with the given ambient pure-action interchanger.
\item[C5, both multiplications.] With $F\to E\to A$, use the common expanded carrier $dR\times_{S^2A}S^2F$. Both routes to the flattened carrier have projections to $dR$ and $SF$ given by identity and $\mu^S_F$. Compatibility with the intermediate $E$ projection is naturality of $\mu^S$. The grade routes agree by Cartesian preservation of accumulation and lift composition.
\item[C6, content unit against context multiplication.] Cartesianness and naturality of $\eta^S$ identify both routes with the same composite display $F\to E\to A$ restricted to the unit body. Its lift/unit grade maps agree by uniqueness.
\item[C7, context associativity.] Nested pullback pasting for three composable displays has the same projections to the original content extension and the final $S$-extension. The grade maps are the coherent composition comparisons for Cartesian lifts.
\item[C8, context right unit.] Removing an identity display before or after pullback gives the identical projection maps and grade lift comparison.
\item[C9, content right unit.] Use $\mu^S(S\eta^S)=\mathrm{id}$ on the representative restricted along the actual right unit body, with the contentad right unit.
\item[C10, context left unit.] Initial identity-display insertion/removal has the same universal projections and lift map on both routes.
\end{description}
In every case (V) identifies the body-side projections, and pullback uniqueness gives the complete context/body equality. Cartesian uniqueness gives the native content-grade equality, even when $d$ is nonfaithful. These arguments are natural in all native arrows because all comparisons were obtained from the same universal factorizations; pure-middle changes use the already verified natural mates. Thus the equations hold in the retained categories, not merely on objects or endpoint images.

Finally, retained comma normalization is proved directly: send a represented pair with bodies $eP\xrightarrow f B\xrightarrow g dQ$ to $gf$; use middle $eP$ for the inverse. On a native pair of maps, apply its content arrow, balance along its context-extension arrow, then apply its context arrow. Balance naturality and composition give the inverse and triangle equations. This does not quotient native arrows. The established pseudomonad composite construction now yields the stated relative multiplication.
\end{proof}

\section{Difference from prior O1 and the generality obtained}
The original independent O1 construction derives occurrence tensor and dependent coefficient transport from a coherent symmetric rig on Set with finitary polynomial occurrences. Here the context doctrine is instead the display-map contextad on an arbitrary pullback category, and the occurrence monad may be any Cartesian monad preserving the display class. The proof uses pullback universal properties rather than finite addressed coefficient calculations. This broadens the base/occurrence assumptions while narrowing the context representation; it does not claim to retain all additional nonfaithful rig coefficients automatically.

The criterion also differs from the older unresolved display/codisplay pullback--pushout proposal. It assumes a valid independent contentad and the stated occurrence display and Cartesian exactness; it derives the mixed law using pullbacks, not an unproved pushout commuting rule. Display-map distributive-law work is prior art, including Tamara von Glehn, \emph{Polynomials, fibrations and distributive laws}, Theory and Applications of Categories 33 (2018), No. 36, pp. 1111--1144, \url{https://www.tac.mta.ca/tac/volumes/33/36/33-36abs.html}, especially its pullback and display-class constructions. This source is also cited in the original dossier. The project-specific comparison and any priority claim require further source-level analysis; no novelty is asserted here.

A shared boundary can be retained by applying the criterion in an appropriate slice base when the independent hypotheses hold there. A display may retain a full joint witness rather than being assembled from independent marginals. But the theorem does not by itself say which physical dependence laws a chosen base encodes, nor that arbitrary boundary forgetting satisfies the hypotheses. Those are semantic questions to test after this general construction, not assumptions hidden in its coherence proof.
\end{document}
